% =====================================================================
%  Contact-aware grasp synthesis and execution
%  Draft section for an ICRA-format paper.
%
%  Notation follows the paper's symbol table (Drake monogram). Symbols
%  introduced here that are NOT yet in that table are collected in
%  Sec.~\ref{sec:new-symbols} at the end of this file so they can be
%  merged in.
%
%  Requires: amsmath, amssymb, bm, algorithm, algpseudocode.
% =====================================================================

\section{Method}
\label{sec:method}

We formulate grasp synthesis as a single nonlinear program that couples
kinematic reachability, contact placement on the object surface, and
wrench feasibility, and we solve it by a sequence of relinearized
subproblems. The resulting contact set is then executed by an internal-force
controller whose commanded squeeze is certified against an explicit
disturbance budget, and whose motion is rate-limited so that the budget
holds by construction.

\subsection{Contact parameterization}
\label{sec:contact-param}

Let the object be represented by a triangulated surface mesh with vertices
$\bm{v}_m \in \mathbb{R}^3$ expressed in $\mathcal{O}_k$. Given a seed contact
$\bm{p}_i^0$ and an outward normal $\bm{n}_i$, we extract the vertices in a
neighbourhood
%
\begin{equation}
  \mathcal{N}_i = \bigl\{ \bm{v}_m : \lVert \bm{P}_i (\bm{v}_m - \bm{p}_i^0) \rVert \le \rho,\;
                   \lvert \bm{n}_i^{\top}(\bm{v}_m - \bm{p}_i^0) \rvert \le \beta_{\mathrm{band}} \bigr\},
\end{equation}
%
where $\bm{P}_i = \bm{I} - \bm{n}_i \bm{n}_i^{\top}$ projects onto the tangent
plane, $\rho$ is the patch radius and $\beta_{\mathrm{band}}$ a thickness
tolerance, and fit both a plane and a quadratic height field to
$\mathcal{N}_i$ in least-squares. Writing the tangential coordinates
$\bm{t}_i = (t_{i,1}, t_{i,2})$, the retained model is
%
\begin{equation}
  h_i(\bm{t}_i) \;=\; -\frac{\kappa_{i,1} t_{i,1}^2 + \kappa_{i,2} t_{i,2}^2}{2 \lVert \nabla \phi_i \rVert},
  \label{eq:patch-height}
\end{equation}
%
with principal curvatures $\kappa_{i,1}, \kappa_{i,2}$ obtained from the
eigendecomposition of the fitted Hessian restricted to the tangent plane, and
$\phi$ the object signed-distance field. The quadratic is accepted only when it
reduces the RMS residual over the planar fit by a factor $\eta$; otherwise
$\kappa_{i,1} = \kappa_{i,2} = 0$ and the patch is treated as locally flat. This
model-selection step matters because scanned meshes carry surface relief that a
curvature estimate will otherwise interpret as genuine geometry.

Each contact then contributes exactly two decision variables,
%
\begin{equation}
  \bm{p}_i(\bm{t}_i) \;=\; \bm{p}_i^0 + t_{i,1}\, \bm{a}_{i,1} + t_{i,2}\, \bm{a}_{i,2} + h_i(\bm{t}_i)\, \bm{n}_i ,
  \label{eq:patch}
\end{equation}
%
where $\bm{a}_{i,1}, \bm{a}_{i,2}$ are the fitted principal axes. Contacts
parameterized this way lie on the surface by construction, so no surface
equality constraint is required and the optimizer cannot leave the object.

\paragraph{Trust region.}
The quadratic \eqref{eq:patch-height} is a local model, and its validity is
bounded by the extent over which it agrees with $\phi$. For each axis and each
direction we march outward until the surrogate departs from the true
signed-distance value by more than $\varepsilon_\phi$, giving asymmetric bounds
%
\begin{equation}
  \underline{t}_{i,d} \le t_{i,d} \le \overline{t}_{i,d}, \qquad d \in \{1,2\},
  \label{eq:trust}
\end{equation}
%
which are tight near an edge and saturate at a user cap on a flat face. Using
one symmetric bound $\min(\lvert \underline{t} \rvert, \overline{t})$ instead
discards usable surface on any patch that is not centred in its face.

\subsection{Grasp synthesis program}
\label{sec:nlp}

Let $\bm{q}^r$ denote the robot configuration and
$\bm{t} = (\bm{t}_1,\dots,\bm{t}_n)$ the stacked contact coordinates. We solve
%
\begin{equation}
\begin{aligned}
  \min_{\bm{q}^r,\, \bm{t},\, \gamma,\, \bm{y}} \quad
    & w_{\mathrm{ik}} J_{\mathrm{ik}} + w_{\mathrm{reg}} J_{\mathrm{reg}}
      + w_{\gamma} J_{\gamma} + w_{y} J_{y} \\[-1pt]
    & \qquad + w_{\alpha} J_{\mathrm{align}} + w_{o} J_{\mathrm{orient}}
      + w_{g} J_{\mathrm{gws}} + w_{s} J_{\mathrm{span}} \\
  \text{s.t.}\quad
    & \bm{q}^r_{\min} \le \bm{q}^r \le \bm{q}^r_{\max}, \\
    & \eqref{eq:trust}, \\
    & \text{wrench feasibility } \eqref{eq:wrench-lp}, \\
    & \mathrm{SDF}(b, o) \ge \epsilon
      \quad \forall\, b \in \mathcal{B}_{\mathrm{collision}} \setminus \mathcal{D},\; o \in \mathcal{O}.
\end{aligned}
\label{eq:nlp}
\end{equation}
%
Each cost term is normalized so that it is of order unity at its reference
level, which lets the weights be read as relative priorities rather than as
unit conversions.

\paragraph{Reachability.}
With $\bm{s}_i(\bm{q}^r)$ the fingertip position at contact $i$ and $r_i$ its
effective pad radius,
%
\begin{equation}
  J_{\mathrm{ik}} = \frac{1}{2 d_{\mathrm{ref}}^2} \sum_{i=1}^{n}
    \bigl\lVert \bm{s}_i(\bm{q}^r) - \bigl(\bm{p}_i(\bm{t}_i) + r_i \bm{n}_i \bigr) \bigr\rVert^2 ,
\end{equation}
%
with $d_{\mathrm{ref}}$ the residual considered acceptable. The target is offset
along the inward normal by the \emph{maximum} pad--mesh offset rather than the
mean, so the converged solution leaves a small positive gap that the squeeze
closes, rather than placing half the contacts in penetration.

\paragraph{Posture.}
$J_{\mathrm{reg}} = \frac{1}{n_{\mathrm{dof}}}
 \lVert \bm{q}^r - \bm{q}_{\mathrm{ref}} \rVert_{\bm{Q}}^2$
ties the solution to the retargeted operator pose, acting as a tie-breaker among
kinematically equivalent grasps and keeping the synthesized configuration close
to the one the operator commanded.

\paragraph{Wrench feasibility.}
We linearize each Coulomb cone with the five generators
%
\begin{equation}
  \bm{f}_{i,j} \in \bigl\{ \bm{0},\;
    (1, \pm\mu_i, 0),\; (1, 0, \pm\mu_i) \bigr\},
  \qquad j = 1,\dots,5,
\end{equation}
%
expressed in $\mathcal{C}_i$, and map them into $\mathcal{O}_k$ as wrench-cone
vertices
%
\begin{equation}
  \bm{F}_{i,j} =
  \begin{bmatrix} \bm{p}_i \times \bm{R}_i \bm{f}_{i,j} \\ \bm{R}_i \bm{f}_{i,j} \end{bmatrix}
  \in \mathbb{R}^6 .
\end{equation}
%
The task disturbance is specified as an acceleration budget rather than a single
wrench: given linear and angular budgets $\bm{a}_{\max}$, $\bm{\alpha}_{\max}$,
object mass $m_k$ and principal inertia $\bm{I}_k$, the induced wrench set is the
box
%
\begin{equation}
  \mathcal{W} = \bigl\{ \bm{w} : \lvert \bm{w}_{\tau} \rvert \le \bm{I}_k \bm{\alpha}_{\max},\;
                   \lvert \bm{w}_{f} \rvert \le m_k \bm{a}_{\max} \bigr\},
  \label{eq:budget}
\end{equation}
%
whose $2^6$ vertices $\bm{w}_\ell$ we enumerate. Gravity enters as a constant
wrench added to every vertex. Feasibility requires that each vertex be
reproducible by cone forces at squeeze $\gamma$:
%
\begin{equation}
\begin{aligned}
  & \bm{V}(1)\, \bm{y}_\ell - \gamma\, \bm{w}_{\mathrm{tan}} = \bm{w}_\ell,
    && \forall \ell, \\
  & \textstyle\sum_{j} y_{\ell,ij} \le \gamma, \quad \bm{y}_\ell \ge \bm{0},
    && \forall \ell, i,
\end{aligned}
\label{eq:wrench-lp}
\end{equation}
%
where the substitution $\bm{y} = \gamma \bm{x}$ renders the condition linear in
$(\bm{y}, \gamma)$, exploiting $\bm{V}(\gamma) = \gamma \bm{V}(1)$. Because the
constraint is embedded in \eqref{eq:nlp} rather than solved as a nested program,
the solver's KKT system couples $\gamma$ directly to $\bm{q}^r$ and $\bm{t}$ and
supplies exact gradients. We penalize $J_\gamma = \gamma / \gamma_{\mathrm{ref}}$
to prefer the least squeeze that certifies the budget, and
$J_y = \sum_\ell \lVert \bm{y}_\ell \rVert^2$ to distribute load across
generators. A per-vertex slack $\bm{s}_\ell$ may be admitted so that a single
geometrically unreachable corner does not render the program infeasible, in
which case feasibility is adjudicated post hoc on $\lVert \bm{s}_\ell \rVert$.

\paragraph{Closure quality.}
Optionally we include the min-weight metric of FRoGGeR~\cite{frogger}. With
$\bm{W} \in \mathbb{R}^{6 \times 5n}$ the matrix of primitive wrenches,
%
\begin{equation}
  J_{\mathrm{gws}} = -\beta^\star, \qquad
  \begin{aligned}
    \beta^\star = \max_{\bm{\alpha}, \beta} \;\; & \beta \\
    \text{s.t.}\;\; & \bm{W}\bm{\alpha} = \bm{0}, \quad
                      \textstyle\sum_j \alpha_j = 1, \\
                    & \bm{\alpha} \ge \beta \bm{1}.
  \end{aligned}
  \label{eq:minweight}
\end{equation}
%
Here $\bm{\alpha}$ is a convex combination of primitive wrenches that cancels to
zero net wrench, and $\beta$ is the smallest weight it requires; $\beta > 0$ iff
the wrench-space origin lies in the interior of the grasp wrench set, and its
magnitude measures how robustly. We deliberately do \emph{not} impose
$\bm{\alpha} \ge \bm{0}$: for a configuration not yet in closure, the correct
value is $\beta < 0$ with some $\alpha_j < 0$, and leaving the sign free keeps
the metric smooth and ascendable from infeasible starts rather than making the
program infeasible. As \eqref{eq:minweight} is embedded in the same program, no
outer loop or subgradient is needed. A companion term
$J_{\mathrm{span}} = -\log\det(\bm{W}\bm{W}^{\top} + \delta \bm{I})$ discourages
rank collapse, which a two-contact pinch approaches whenever the contacts become
collinear.

\paragraph{Geometric conditioning.}
$J_{\mathrm{align}} = \lVert \hat{\bm{g}} - \bm{n}_1 \rVert^2$ with
$\hat{\bm{g}} = (\bm{p}_2 - \bm{p}_1)/\lVert \bm{p}_2 - \bm{p}_1 \rVert$ rewards
antipodal opposition, and $J_{\mathrm{orient}}$ penalizes deviation between each
pad axis and its inward contact normal so the pad meets the surface flush.

\paragraph{Collision.}
Robot--object and robot--environment clearances use a smooth (softplus)
sphere--primitive distance so the constraint is differentiable. The fingertips
$\mathcal{D}$ are excluded from the object clearance constraint, since they are
required to touch, but retain the support-surface constraint so they cannot be
driven through it.

\subsection{Relinearization}
\label{sec:picard}

Contact normals and the frames $\bm{R}_i$ are held fixed within a solve, which
keeps \eqref{eq:nlp} smooth. We therefore iterate: solve, re-anchor each seed at
the returned contact, refit \eqref{eq:patch-height}, and repeat for a fixed
number of stages. Convergence is assessed by whether contacts remain interior to
\eqref{eq:trust}; a solution resting on a trust-region bound indicates that the
local model, not the objective, is limiting the step.

\subsection{Execution}
\label{sec:execution}

The synthesized contacts are executed in five phases: kinematic approach, settle,
squeeze, transport, and release.

\paragraph{Certified squeeze.}
The squeeze commanded at run time is \emph{not} the program's internal $\gamma$.
Once the hand is in contact we re-measure the achieved contact geometry
$(\bm{p}_i, \bm{R}_i)$ and re-solve \eqref{eq:wrench-lp} as a linear program
over the realized configuration, obtaining $\gamma^\star$, the least squeeze
that renders every vertex of \eqref{eq:budget} resistible. Under zero
disturbance this reduces to $\gamma^\star = m_k g / (2\mu)$ for an antipodal
pinch, i.e.\ exactly the friction required to carry the object's weight.
Infeasibility of any vertex is reported rather than absorbed.

\paragraph{Budget enforcement.}
The certificate $\gamma^\star$ is meaningful only if the executed motion remains
inside \eqref{eq:budget}. We therefore rate-limit the commanded end-effector
twist so that
%
\begin{equation}
  \bm{v}^{c}_{k+1} = \bm{v}^{c}_{k} +
    \mathrm{clip}\bigl( \bm{v}^{d}_{k} - \bm{v}^{c}_{k},\; \pm \bm{a}_{\max} \Delta t \bigr),
  \label{eq:slew}
\end{equation}
%
with $\bm{v}^d$ the desired and $\bm{v}^c$ the commanded twist. The commanded
acceleration is then bounded by $\bm{a}_{\max}$ by construction, so the budget
the squeeze was certified against is enforced rather than assumed. This is what
permits the traverse speed to be raised without renegotiating the grasp
force: peak speed sets the duration of the transfer, while \eqref{eq:slew} sets
what the grasp must survive. Joint rates follow from a damped least-squares
resolved-rate law with damping raised near kinematic singularities.

\paragraph{Phase-dependent compliance.}
Closing and holding impose opposing requirements on finger impedance. During
closing a compliant finger lets the internal-force term dominate and seat the
contacts; under load the same compliance is back-driven and bleeds normal force
until the pinch fails. We therefore switch the finger gains at the transition
from squeeze to transport.

% =====================================================================
\subsection*{Symbols to add to the notation table}
\label{sec:new-symbols}
% =====================================================================
% Suggested rows, following the existing Drake-monogram convention.
%
% --- Contacts -------------------------------------------------------
%   $\bm{t}_i \in \mathbb{R}^2$        Tangential patch coordinates of contact $i$
%   $\bm{a}_{i,1}, \bm{a}_{i,2}$       Principal axes of the fitted patch at contact $i$
%   $\kappa_{i,1}, \kappa_{i,2}$       Principal curvatures of the fitted patch
%   $h_i(\bm{t}_i)$                    Local surface height field, Eq. (2)
%   $\bm{n}_i \in \mathbb{S}^2$        Outward surface normal at contact $i$
%   $r_i$                              Effective fingertip pad radius at contact $i$
%   $\phi(\cdot)$                      Object signed-distance field
%   $\rho, \beta_{\mathrm{band}}$      Patch fit radius / band thickness
%   $\eta$                             Plane-vs-quadratic RMS gain threshold
%   $\varepsilon_\phi$                 Surrogate--SDF agreement tolerance
%   $\underline{t}_{i,d}, \overline{t}_{i,d}$  Asymmetric trust-region bounds
%
% --- Task / feasibility ---------------------------------------------
%   $\bm{a}_{\max}, \bm{\alpha}_{\max}$  Linear / angular acceleration budget
%   $m_k, \bm{I}_k$                    Mass and principal inertia of object $k$
%   $\mathcal{W}$                      Task wrench box, Eq. (9)
%   $\bm{w}_\ell$                      Vertex $\ell$ of $\mathcal{W}$
%   $\bm{s}_\ell \in \mathbb{R}^6$     Per-vertex wrench slack
%   $\gamma^\star$                     Certified squeeze from the realized grasp
%   $\bm{W} \in \mathbb{R}^{6\times 5n}$  Primitive wrench matrix
%   $\bm{\alpha}, \beta$               Min-weight LP variables, Eq. (11)
%   $\beta^\star$                      Min-weight metric value
%   $\delta$                           Span-term regularizer
%
% --- Costs ----------------------------------------------------------
%   $J_{(\cdot)}$                      Normalized cost terms
%   $w_{(\cdot)}$                      Corresponding weights
%   $d_{\mathrm{ref}}$                 Reference IK residual
%   $\gamma_{\mathrm{ref}}$            Reference squeeze (task load)
%
% --- Execution ------------------------------------------------------
%   $\bm{v}^{d}, \bm{v}^{c}$           Desired / commanded end-effector twist
%   $\Delta t$                         Control period
%
% NOTE ON A COLLISION: the table already uses $\bm{\tau}_i$ for the tangential
% internal force at contact $i$, and $\gamma$ for the squeeze factor. Above,
% $\bm{w}_\tau$ denotes the torque block of a wrench -- consider renaming that
% to $\bm{w}_{m}$ (moment) to avoid overloading $\tau$. Likewise the table's
% $\mathcal{O}$ is the set of scene objects while $\mathcal{O}_k$ is a body
% frame; that is workable but worth a footnote.
